첫 번째 줄에 지출 내역의 개수 $N$이 주어진다. $(5 \le N \le 1\,000)$

다음 $N$개의 줄에 걸쳐 각 지출 내역이 $\texttt{"MM DD CATEGORY AMOUNT"}$ 형식으로 주어진다.

\begin{itemize}
\item $\texttt{MM}$과 $\texttt{DD}$는 각각 월과 일을 의미하며, 항상 두 자리 숫자로 주어진다. 한 자리 숫자는 앞에 0을 붙여 표기한다. $(01 \le \texttt{MM} \le 12, 01 \le \texttt{DD} \le 31)$ 
\item $\texttt{CATEGORY}$는 \texttt{Food, Hobbies, Living, Healthcare, Transportation} 중 하나이다.
\item $\texttt{AMOUNT}$는 $1$ 이상 $1\,000\,000$ 이하이다.
\item 주어지는 모든 지출 내역의 월은 동일하며, 모든 날짜는 실제 달력에 존재하는 날이다. 단, 윤일($2$월 $29$일)은 주어지지 않는다.
\item 지출 내역은 날짜의 오름차순으로 주어진다. 같은 날짜의 지출 내역이 여러 개 주어질 수 있다.
\item 모든 카테고리는 적어도 한 번 이상 등장한다.
\item 총지출 금액이 최대인 카테고리와 최소인 카테고리는 각각 유일하다.
\end{itemize}

